% Canonical equation -- KH instability linear stability dispersion (KH-LS)
% kh-sim canonical-equations / v0.1
% Companion to kh-linear-stability.md
% Bibliography: ../reference-bibliography/refs.bib
% Licence: CC-BY-SA-4.0

\documentclass[11pt,a4paper]{article}
\usepackage[utf8]{inputenc}
\usepackage[T1]{fontenc}
\usepackage{amsmath, amssymb, amsthm}
\usepackage{geometry}
\usepackage{booktabs}
\usepackage{hyperref}
\usepackage[backend=biber,style=numeric]{biblatex}
\addbibresource{../reference-bibliography/refs.bib}

\geometry{margin=2.4cm}

\title{Canonical equations:\\ Kelvin--Helmholtz linear stability dispersion}
\author{kh-sim project}
\date{v0.1 -- 2026-05-12}

\theoremstyle{definition}
\newtheorem{definition}{Definition}
\theoremstyle{plain}
\newtheorem{proposition}{Proposition}

\begin{document}
\maketitle

This document states the linear-stability dispersion relation
$\sigma(k_x)$ for the kh-sim canonical initial conditions
(tanh shear profile of Eq.\ KH-GE-4). It is the analytical anchor for
one of the property tests in the phase-7 cross-backend property suite:
\emph{the simulated early-time growth rate of a small-amplitude
perturbation at mode $k_x$ must match the analytic $\sigma(k_x)$
within tolerance.}

This is the most demanding property test in the suite -- it certifies
that the implementation correctly solves the linearised Navier--Stokes
equations, not just that it conserves a few invariants.

\section{The Rayleigh stability equation (Eq.\ KH-LS-1)}

\begin{definition}[Inviscid linear stability for a parallel shear flow]
Consider the inviscid limit ($\mathrm{Re} \to \infty$) of the kh-sim
governing equations (KH-GE-1 with $\nu = 0$). Linearise about a
parallel base flow $U(y)$ and seek perturbation normal modes of the
form
\[
\omega'(x, y, t) = \tilde\omega(y)\,e^{i k_x x + \sigma t},
\qquad
\psi'(x, y, t) = \tilde\psi(y)\,e^{i k_x x + \sigma t}.
\]
The resulting eigenvalue problem is the \emph{Rayleigh stability
equation}~\cite{DrazinReid2004}:
\begin{equation}
\label{eq:khls1}
(U - c)\left(\frac{d^2 \tilde\psi}{dy^2} - k_x^2\,\tilde\psi\right)
- U''\,\tilde\psi = 0,
\tag{KH-LS-1}
\end{equation}
where $c = -i\sigma / k_x$ is the complex phase speed (real part:
propagation speed; imaginary part: growth rate divided by $k_x$).
\end{definition}

\section{tanh-profile dispersion (Eq.\ KH-LS-2)}

\begin{proposition}[tanh-profile inviscid dispersion relation, Drazin \& Reid \S4.5]
For the base flow
\[
U(y) = U_0\,\tanh\!\left(\frac{y - L_y/2}{\delta}\right),
\]
the Rayleigh equation \eqref{eq:khls1} admits a known analytic
solution. The growth rate $\sigma(k_x)$ is real and positive
(\emph{unstable}) for $k_x \delta < 1$, and purely imaginary
(neutrally stable oscillation) for $k_x \delta > 1$. The
fastest-growing mode satisfies
\begin{equation}
\label{eq:khls2}
k_{x, \max}\,\delta \approx 0.4446,
\qquad
\sigma_{\max} \approx 0.19\,\frac{U_0}{\delta}.
\tag{KH-LS-2}
\end{equation}
\end{proposition}

The full $\sigma(k_x)$ curve over $0 < k_x \delta < 1$ is most
conveniently obtained by direct numerical solution of \eqref{eq:khls1}
(shooting method or boundary-value solver on a truncated $y$-domain).
See Schmid \& Henningson~\cite{SchmidHenningson2001} \S 3 for tabulated
benchmark values.

\section{The kh-sim canonical mode (Eq.\ KH-LS-3)}

The kh-sim canonical IC (Eq.\ KH-GE-4) excites perturbation mode
$k = 2$ over domain length $L_x = 1.0$, giving physical wavenumber
\[
k_x = \frac{2\pi k}{L_x} = 4\pi \approx 12.566\;\text{rad/m}.
\]
With $\delta = 0.05\,L_y = 0.025\;\text{m}$ and $U_0 = 1\;\text{m/s}$,
the non-dimensional combination is
\[
k_x\,\delta = 12.566 \cdot 0.025 \approx 0.314,
\]
comfortably below the marginal-stability cutoff $k_x \delta = 1$.
The mode is in the unstable band. Reading the tanh-profile dispersion
curve at $k_x \delta = 0.314$ gives
\begin{equation}
\label{eq:khls3}
\sigma_{\text{expected}} \approx 0.18\,\frac{U_0}{\delta}
\approx 7.2\;\text{s}^{-1}.
\tag{KH-LS-3}
\end{equation}
(Provisional numerical value -- to be confirmed in phase 7 with an
explicit numerical solve of the Rayleigh equation~\eqref{eq:khls1}
for the kh-sim default parameters.)

\section{Property test specification (phase 7)}

\begin{definition}[Linear-growth-rate property test]
Run kh-sim with the canonical IC of Eq.\ KH-GE-4 except with a
reduced perturbation amplitude $A = 10^{-4}$ (instead of the canonical
$A = 10^{-2}$), to ensure the perturbation remains in the linear
regime during the test interval. Measure the kinetic energy
\[
E_k(t) = \tfrac{1}{2}\,\mathbb{E}\!\left[ u^2 + v^2 \right]
\]
at output times $t \in \{0, t_1, t_2, \ldots\}$ with $t_n A \ll 1$ so
the simulation is still linear. The simulated growth rate is
\[
\sigma_{\text{simulated}} = \tfrac{1}{2}\,\frac{d\ln E_k(t)}{dt}\bigg|_{t \to 0^+}
\]
(the factor of $1/2$ converts from energy growth rate $\propto e^{2\sigma t}$
back to amplitude growth $\propto e^{\sigma t}$).
\end{definition}

\begin{proposition}[Property-test gate]
For kh-sim at the canonical IC parameters with $A$ reduced to
$10^{-4}$, $\sigma_{\text{simulated}}$ shall match
$\sigma_{\text{expected}}$ from \eqref{eq:khls3} within:
\begin{itemize}
\item $\lvert \Delta\sigma/\sigma \rvert < 5\%$ at canonical
      resolution $(N_x, N_y) = (128, 64)$ with default $\Delta t = 10^{-3}$;
\item $\lvert \Delta\sigma/\sigma \rvert < 1\%$ at well-resolved
      settings $(N_x, N_y) = (256, 128)$ with $\Delta t \leq 5 \times 10^{-4}$.
\end{itemize}
\end{proposition}

This is the canonical physics-correctness property test for kh-sim.
It anchors not only numerical-method correctness (which the
TC-NUM-KH-001..008 series in
\texttt{kh-numerical-methods.tex} already covers) but also that the
implementation correctly recovers a non-trivial \emph{physical}
prediction from the equations.

\section{Outside scope of this document}

The following are deliberately excluded from this canonical statement
and treated in companion or future material:
\begin{itemize}
\item Viscous corrections to $\sigma(k_x)$ for finite~$\mathrm{Re}$.
      For $\mathrm{Re} = 1000$ (kh-sim default), the correction is
      small but non-zero. Covered in
      Drazin \& Reid~\cite{DrazinReid2004} \S 4.6 and Schmid \&
      Henningson~\cite{SchmidHenningson2001} \S 3.4.
\item Non-linear billow roll-up dynamics. The linear theory captures
      only the initial exponential growth; the visible
      ``rolling vortex structure'' of the Kelvin--Helmholtz
      instability is a non-linear saturation phenomenon (vortex sheet
      roll-up, vortex pairing).
\item Density-stratified Kelvin--Helmholtz (Chandrasekhar
      \cite{Chandrasekhar1981}; Schmid \& Henningson \S 3.5).
      kh-sim assumes uniform density.
\end{itemize}

\printbibliography

\end{document}
